\begin{description}
\item[(a)] The distance traveled by the gas particles of inner and outer radius
  can be calculated as follows:
  \begin{align*}
    S_{\mathrm{out}} &= d_\oplus(\theta'_{\mathrm{out}} - \theta_{\mathrm{out}})
      = 100\,\mathrm{pc}\cdot(8.0' - 4.0') = 24\times10^{3}\ \mathrm{au} \\
    S_{\mathrm{in}} &= d_\oplus(\theta'_{\mathrm{in}} - \theta_{\mathrm{in}})
      = 100\,\mathrm{pc}\cdot(7.0' - 3.5') = 21\times10^{3}\ \mathrm{au}
  \end{align*}
  \hfill(2.0 points)
  \begin{align*}
    \therefore\; v_{max} &= \frac{S_{\mathrm{out}}}{t_0} = 57.0\ \mathrm{km/s} \\
    v_{min} &= \frac{S_{\mathrm{in}}}{t_0} = 50.0\ \mathrm{km/s}
  \end{align*}
  \hfill(2.0 points)
\item[(b)] At the inner edge, escape velocity from the white dwarf can be
  maximal ($1.44M_\odot$ -- Chandrasekhar limit)
  \[
    v_{esc} = \sqrt{\frac{2G\cdot1.44M_\odot}{S_{\mathrm{in}}}}
      \approx 0.35\ \mathrm{km/s}
  \]
  \hfill(2.0 points)

  Since $v_{esc} \ll v_{min}$, the no-gravity scenario can be applied. YES
  \hfill(1.0 points)
\item[(c)]
  \begin{align*}
    \phi &= 8' - 7.5' = 0.5' = 30'' \ge \frac{1.22\lambda}{D} \\
    \therefore\; D &\ge \frac{1.22\lambda}{\phi}
      = \frac{1.22\times5\times10^{-7}\times206265}{30} = 4.19\ \mathrm{mm}
  \end{align*}
  \hfill(2.0 points)

  Hence YES
  \hfill(1.0 points)
\end{description}
